\documentclass[11pt,a4paper]{article}
\usepackage[T1]{fontenc}
\usepackage{lmodern,amsmath,amssymb,amsthm,mathtools,microtype,xcolor}
\usepackage[margin=26mm]{geometry}
\usepackage[colorlinks=true,linkcolor=blue!45!black,citecolor=blue!45!black,urlcolor=blue!45!black]{hyperref}
\newtheorem{theorem}{Theorem}[section]
\newtheorem{lemma}[theorem]{Lemma}
\newtheorem{proposition}[theorem]{Proposition}
\newtheorem{corollary}[theorem]{Corollary}
\theoremstyle{remark}\newtheorem{remark}[theorem]{Remark}
\newcommand{\Q}{\mathbb Q}\newcommand{\Z}{\mathbb Z}
\setlength{\emergencystretch}{2em}
\title{Separation and denominator conditions in quadratic prime patterns}
\author{}\date{9 November 2026}
\begin{document}\maketitle
\begin{abstract}
We derive a separation property for quadratic prime patterns from the
weighted estimates in Sun's simultaneous-avoidance argument. The selected
primes can avoid every selected opposite elliptic reduction order while
retaining fixed congruences and arbitrarily large prescribed denominator
divisibility. We also give a denominator normalization that permits all
local weights of order $m$ to be replaced by powers of one weight. These
consequences use the cited analytic estimates and elementary valuation
arguments; they do not require a further prime-distribution hypothesis.
\end{abstract}

\section{The precise analytic input}
Fix distinct integers $c_1,\ldots,c_k$, with $k\ge2$, and put
\[
 p_j=u+c_jw^2.
\]
The fixed congruences, Gaussian ray conditions and reduction-order
avoidance conditions in \cite[Theorem 4.1]{Sun} are imposed throughout.
A fixed dilation of $w$ can be absorbed into the integers $c_j$.
All such data are chosen before the parameter $X$ tends to infinity.

We use the following quantitative content of
\cite[Propositions 5.1 and 6.2, Section 6.1]{Sun}. There are constants
$c,C>0$ and a sufficiently slowly decreasing function $\tau(X)>0$,
with $\tau(X)\to0$, such that, for
\[
 T=\lfloor\tau(X)\sqrt X\rfloor,
\]
the admissible pairs in $X/2\le u\le X$, $1\le w\le T$ carry
nonnegative weights $\omega_X(u,w)$ satisfying
\begin{equation}\label{eq:count}
 \sum_{u,w}\omega_X(u,w)\ge cXT,
 \qquad 0\le\omega_X(u,w)\le C(\log X)^k.
\end{equation}
The weights vanish outside admissible pairs. The lower bound is after
the removal of all self-order and cross-order obstructions, not merely
after a finite list of congruences.

For clarity, the order of limits in the cited argument is essential.
Its finite cutoff, denoted by $L$ there, is fixed first. The surviving
weight is at least
\[
 \frac{c_1XT}{8(\log L)^{4k}}
\]
for a positive constant $c_1$ and all sufficiently large $X$.
This gives the first inequality in \eqref{eq:count} with a fixed
positive constant. Each surviving weight is a product of $k$ prime
logarithms, up to fixed normalizations, giving the second inequality.
The diagonal construction of $\tau$ allows additional lower thresholds
which hold eventually for every fixed positive value of $\tau$.

All assertions below that select prime patterns depend on these cited
estimates. The elementary statements about weights, integers and
valuations are proved here.

\section{Removing a thin strip}
\begin{lemma}\label{lem:strip}
Assume \eqref{eq:count}. Let $0<\alpha<1/2$ be fixed.
The function $\tau$ may be chosen sufficiently slowly that admissible
pairs with $w\ge X^\alpha$ still have total weight at least $cXT/2$
for every sufficiently large $X$.
\end{lemma}
\begin{proof}
The strip $w<X^\alpha$ contains at most $X^{1+\alpha}+X$
integer pairs. Its weight is at most
\[
 C(X^{1+\alpha}+X)(\log X)^k.
\]
Choose the diagonal thresholds so that
\[
 \frac{\tau(X)X^{1/2-\alpha}}{(\log X)^k}\longrightarrow\infty.
\]
This is compatible with every prior threshold: for a fixed positive
value $\tau_j$, the displayed quotient tends to infinity, so the start
of the $j$th diagonal interval can be increased to meet this additional
condition. Also require $\tau(X)\sqrt X\ge2$, so
$T\ge\tau(X)\sqrt X/2$. The strip weight divided by $XT$ tends
to zero. Subtracting it from \eqref{eq:count} proves the assertion.
\end{proof}

\begin{theorem}\label{thm:separation}
For each selected prime $p_j$, let $d_j$ be a positive integer such that
\[
 d_j=p_j+O(\sqrt{p_j}),\qquad 4\mid d_j,
\]
where the implied constant is uniform in $j$ and the pattern.
Under the cited analytic input, the pattern can be chosen so that
\[
                       p_i\nmid d_j\qquad(1\le i,j\le k).
\]
The fixed congruence and avoidance conditions remain satisfied, all
selected primes tend to infinity, and $u/w^2$ tends to infinity.
\end{theorem}
\begin{proof}
Apply Lemma~\ref{lem:strip} with $1/4<\alpha<1/2$.
Since $w^2\le T^2=o(X)$, every $p_j/u$ tends to one, uniformly
over the selected patterns. The same is true of $d_j/u$.
Consequently $1/2<d_j/p_i<2$ for sufficiently large $X$.
Divisibility would therefore imply $d_j=p_i$.

For $i\ne j$, let $c_*:=\min_{i\ne j}|c_i-c_j|>0$. Then
\[
 |p_i-p_j|\ge c_*w^2\ge c_*X^{2\alpha},
\]
whereas $|d_j-p_j|=O(\sqrt X)$. Because $2\alpha>1/2$,
equality is impossible for sufficiently large $X$.
For $i=j$, the divisibility $4\mid d_i$ excludes $d_i=p_i$,
since $p_i$ is an odd prime. The other conditions follow from the
support of the weights and from $u\ge X/2$, $w^2=o(X)$.
\end{proof}

For the elliptic curves used in the defining-polynomial constructions,
Hasse's bound supplies the displayed estimate for $d_j$, and the fixed
Gaussian congruences make the opposite order divisible by four.
Thus Theorem~\ref{thm:separation} gives precisely the separation used
to impose additional formal depth on the primary points.
The application of this observation to those constructions is already
present in \cite[Version 9 separation subsection]{Upstream}; the
argument above states its quantitative source and permitted limits.

\section{Simultaneous denominator preparation}
Let $S$ be a finite set of rational primes, and let
$a_1,\ldots,a_s\in\Q^\times$ be fixed.

\begin{lemma}\label{lem:denominator}
There is a positive integer $D$ such that, whenever $u$ is a prime
outside $S$ and $w$ is a positive integer, the numbers
\[
 \lambda=\frac{u}{D^2w^2},\qquad t_i=a_i\lambda,
 \qquad\delta_i=t_i-1
\]
satisfy $v_p(\delta_i)=v_p(t_i)<0$ for every $p\in S$ and every $i$.
The integer $D$ can also satisfy any prescribed finite set of
divisibility conditions.
\end{lemma}
\begin{proof}
Choose $D$ with $2v_p(D)>\max_i v_p(a_i)$ for every $p\in S$,
and enlarge it to satisfy the other fixed divisibilities. Since
$v_p(u)=0$ and $v_p(w)\ge0$,
\[
 v_p(t_i)=v_p(a_i)-2v_p(D)-2v_p(w)<0.
\]
The two terms of $t_i-1$ have distinct valuations, so its valuation
equals $v_p(t_i)$. All requirements concern a finite set of primes.
\end{proof}

This choice can be made after fixing the scales and the finite
preparation set, but before selecting primes. Multiplying the
denominator dilation only changes fixed quadratic coefficients of the
contact forms. The coefficient of the zero contact is zero, so $u$
itself is one of the selected primes; the lower cutoff keeps it outside
$S$. This is within the fixed-dilation scope of the cited theorem.

\section{Powers of one local weight}
For nonzero rational $R,\delta$, define
\[
 q_m=R^{1-2m}\delta^{-2m-2},\qquad
 q=q_1=R^{-1}\delta^{-4},\qquad m\ge1.
\]
\begin{proposition}\label{prop:powers}
The alternative weights $q'_m=q^{2m-1}$ satisfy
\[
                 q'_m=q_m\delta^{-6(m-1)}.
\]
At every valuation where $\delta$ is a unit, $q'_m$ and $q_m$ have
the same valuation. At every valuation where $v(\delta)\le0$,
$v(q'_m)\ge v(q_m)$. In particular, a lower bound for
$v(q_mD^2)$ remains true with $q'_m$ in place of $q_m$.
\end{proposition}
\begin{proof}
The exponent of $R$ in $q^{2m-1}$ is $1-2m$ and its exponent of
$\delta$ is $4-8m$. Subtracting the exponent $-2m-2$ in $q_m$
gives $-6(m-1)$. Taking valuations proves the assertions.
\end{proof}

\begin{corollary}\label{cor:preservation}
Consider factors
\[
 1+\sum_{j=0}^{r-1}q_{m_j}^{2^j}D_j^{2^{j+1}}
\]
in a local test. At contact places where $v(R)$ is odd and $\delta$
is a unit, replacing $q_{m_j}$ by $q^{2m_j-1}$ preserves every
valuation and the distinct-valuation argument for its summands.
At the finite places prepared by Lemma~\ref{lem:denominator}, any
prescribed positive lower bounds on the summand valuations are
preserved. If, as $\lambda\to+\infty$, one has
$R\asymp\lambda^{-1}$, $\delta\asymp\lambda$, and all $D_j$
bounded at the real embeddings, the modified factors tend to one.
\end{corollary}
\begin{proof}
The finite-place conclusions follow term by term from
Proposition~\ref{prop:powers}. At a contact place, the valuation of
the $j$th nonconstant summand is $2^j$ times an odd integer; these
valuations and zero are pairwise distinct. At infinity,
$q=O(\lambda^{-3})$, hence
$q^{2m_j-1}=O(\lambda^{-3(2m_j-1)})$. Each nonconstant summand
therefore tends to zero.
\end{proof}

The corollary is a statement about the specified local tests and their
simultaneous preparation. A complete defining-polynomial construction
must still prove the pole realization and index-height steps, specify
its equations, and check their possible zero strata. Those separate
arithmetic requirements are unaffected by this normalization.

\begin{thebibliography}{9}
\bibitem{Sun} Z.-W. Sun, \emph{86 unknowns for Hilbert's tenth problem
over the rationals}, arXiv:2610.20691v2 (2026), Theorem 4.1,
Propositions 5.1 and 6.2, Sections 6.1--8.3.
\url{https://arxiv.org/html/2610.20691v2}.
\bibitem{Upstream} \emph{$\Z$ is Diophantine in $\Q$}, Version 13,
7 November 2026, retained Version 9 subsection on separation and
asymmetric formal depth. Commit
\texttt{5bbd74f6a479776a332cbe036809d5dddae63720}.
\url{https://github.com/elliotglazer/diophantine-z-q}.
\bibitem{Silverman} J. H. Silverman, \emph{The Arithmetic of Elliptic
Curves}, second edition, Graduate Texts in Mathematics 106,
Springer, 2009, Chapter V, Theorem 1.1.
\end{thebibliography}
\end{document}
